% Glossar: Begriffe einmal hier definieren.
% \begriff{schluessel}{Name}{Kurzfassung für die Fußnote}{ausführliche Erklärung im Glossar}
% Im Text: \fb{schluessel} -> Name mit Link ins Glossar, beim ersten Auftreten zusätzlich Fußnote mit der Kurzfassung.
\begriff{recall}{Recall}
  {Anteil der relevanten Einträge, die ein Verfahren findet. Ein Recall von 0,69 bedeutet: Von 100 relevanten Absätzen wurden 69 gefunden und 31 übersehen.}
  {Anteil der tatsächlich relevanten Einträge, die ein Verfahren findet. Maßstab ist eine Referenzliste, die von Fachleuten erstellt wurde (Goldstandard). Ein Recall von 0,69 bedeutet: Von 100 relevanten Absätzen hat das Verfahren 69 gefunden und 31 übersehen. Wie viele nicht relevante Absätze das Verfahren zusätzlich meldet, misst der Recall nicht; das beschreibt die Präzision.}
\begriff{praezision}{Präzision}
  {Anteil der gemeldeten Einträge, die tatsächlich relevant sind. Eine Präzision von 0,64 bedeutet: Von 100 gemeldeten Absätzen sind 64 relevant.}
  {Anteil der gemeldeten Einträge, die tatsächlich relevant sind. Eine Präzision von 0,64 bedeutet: Von 100 gemeldeten Absätzen sind 64 relevant und 36 überflüssig. Ein Verfahren, das alles meldet, erreicht einen hohen Recall, aber eine niedrige Präzision.}
